% !TeX program = lualatex
% ================================================================
% Urdu Parallel Text Version of SequencesNthTermIdeasStarters
%
% Generated by the server using the following settings:
%
%   language            Urdu (urd)
%   text direction      right to left
%   font                Noto Nastaliq Urdu
%   fallback font       Noto Serif
%   built from          latex/starters/targeted/KS3/sequences/sequencesNthTermIdeasStarters.tex
%   vocabulary key      latex/starters/targeted/KS3/sequences/sequencesNthTermIdeasStarters/L2/urd/sequencesNthTermIdeasStarters_urduKey.tex
%   tier 3 vocabulary   green   RGB 0 128 0
%   English text        black   RGB 0 0 0
%   Urdu text           purple  RGB 112 48 160
%   layout macros       version 3
%
% Please don't edit any information in this comment block - the server
% compares it against the current settings and rebuilds this file when
% they differ, so changes here would be overwritten. However, feel free
% to make updates to the LaTeX below if you want to edit it.
% ================================================================

\documentclass{beamer}
% ================================================================
% Copied from the English sheet this was translated from, so that
% anything it sets up for itself still works here
% ================================================================
\usepackage{tikz}
\usepackage{amsmath}

% ---- Minimal styling ----
\setbeamertemplate{navigation symbols}{}
\setbeamertemplate{footline}{}
\setbeamercolor{background canvas}{bg=white}
\setbeamercolor{frametitle}{fg=black}
\setbeamerfont{frametitle}{size=\Large, series=\bfseries}
\usefonttheme{professionalfonts}

% ---- Global tikz baseline ----
\tikzset{every picture/.style={baseline=(current bounding box.south)}}

% ================================================================
% TikZ pattern commands
% ================================================================

% Matchstick squares: n squares in a row
\newcommand{\matchsquares}[2][black]{%
  \begin{tikzpicture}[scale=0.60]
    \pgfmathsetmacro\n{#2}
    \color{#1}
    \foreach \x in {0,...,\n} {
      \draw[very thick, cap=round] (\x,0) -- (\x,1);
    }
    \foreach \x in {0,...,\the\numexpr\n-1\relax} {
      \draw[very thick, cap=round] (\x,0)   -- (\x+1,0);
      \draw[very thick, cap=round] (\x,1)   -- (\x+1,1);
    }
  \end{tikzpicture}%
}

% Matchstick triangles: chain of n triangles
\newcommand{\matchtriangles}[2][black]{%
  \begin{tikzpicture}[scale=0.60]
    \pgfmathsetmacro\n{#2}
    \pgfmathsetmacro\h{0.866}
    \color{#1}
    \draw[very thick] (0,0) -- (\n,0);
    \draw[very thick] (0,0)
      \foreach \k in {1,...,\n} {
        -- ({\k - 0.5}, {\h}) -- (\k,0)
      };
  \end{tikzpicture}%
}

% Rectangular dots: n columns x (n+1) rows
\newcommand{\rectdots}[2][black]{%
  \begin{tikzpicture}[scale=0.60]
    \pgfmathsetmacro\n{#2}
    \color{#1}
    \foreach \i in {0,...,\the\numexpr\n-1\relax} {
      \foreach \j in {0,...,\n} {
        \fill (0.4*\i, 0.4*\j) circle (0.07);
      }
    }
  \end{tikzpicture}%
}

% Cross / plus dots
\newcommand{\crossdots}[2][black]{%
  \begin{tikzpicture}[scale=0.60]
    \pgfmathsetmacro\n{#2}
    \color{#1}
    \fill (0,0) circle (0.07);
    \ifnum\n>1
      \foreach \d in {1,...,\the\numexpr\n-1\relax} {
        \fill (0,   0.4*\d) circle (0.07);
        \fill (0,  -0.4*\d) circle (0.07);
        \fill ( 0.4*\d, 0)  circle (0.07);
        \fill (-0.4*\d, 0)  circle (0.07);
      }
    \fi
  \end{tikzpicture}%
}

% 4-times-table dots: n rows of 4 dots
\newcommand{\fourtabledots}[2][black]{%
  \begin{tikzpicture}[scale=0.60]
    \pgfmathsetmacro\n{#2}
    \color{#1}
    \foreach \row in {0,...,\the\numexpr\n-1\relax} {
      \foreach \col in {0,1,2,3} {
        \fill (0.4*\col, 0.4*\row) circle (0.07);
      }
    }
  \end{tikzpicture}%
}

% ================================================================
% Answer-overlay helpers
% ================================================================
\newcommand{\ablank}[1]{%
  \alt<2>{\textcolor{red}{#1}}{\underline{\phantom{#1}}}%
}
\newcommand{\ashow}[1]{\uncover<2->{\textcolor{red}{\small #1}}}
\newcommand{\ashowq}[1]{\alt<2->{\textcolor{red}{\small #1}}{?}}

% Answers drawn straight on to a TikZ diagram, revealed on overlay 2 like \ashow.
% Space is reserved on overlay 1, so the diagram does not move between slides.
\usetikzlibrary{overlay-beamer-styles}
\tikzset{
  ans/.style     = {red, thick, visible on=<2->},
  ansfill/.style = {red!15, visible on=<2->},
  anslab/.style  = {red, font=\tiny, visible on=<2->},
}

% ================================================================

% Characters this language's font has not got are borrowed from another,
% rather than being silently dropped - see docs/translated-sheet-latex.md
\directlua{%
  if luaotfload and luaotfload.add_fallback then
    luaotfload.add_fallback("eallatinfallback",
      { "Noto Serif:mode=harf;" })
  else
    tex.error("this luaotfload is too old to register a fallback font")
  end
}
\usepackage[english,bidi=basic]{babel}
\babelprovide[import]{urdu}
\babelfont[urdu]{rm}[Renderer=Harfbuzz, BoldFont={Noto Nastaliq Urdu}, ItalicFont={Noto Nastaliq Urdu}, BoldItalicFont={Noto Nastaliq Urdu}, RawFeature={fallback=eallatinfallback}]{Noto Nastaliq Urdu}
\babelfont[urdu]{sf}[Renderer=Harfbuzz, BoldFont={Noto Nastaliq Urdu}, ItalicFont={Noto Nastaliq Urdu}, BoldItalicFont={Noto Nastaliq Urdu}, RawFeature={fallback=eallatinfallback}]{Noto Nastaliq Urdu}
\newcommand{\ealtext}[1]{\foreignlanguage{urdu}{#1}}
\newcommand{\ealtextblock}[1]{{\raggedleft\foreignlanguage{urdu}{#1}\par}}
% ================================================================
% EAL layout helpers
% ================================================================
\usepackage{xcolor}

\definecolor{ealkeycolour}{RGB}{0,128,0}
\definecolor{ealenglish}{RGB}{0,0,0}
\definecolor{ealtwocolour}{RGB}{112,48,160}

\newsavebox{\ealboxen}
\newsavebox{\ealboxtr}
\newlength{\ealglosswd}
\newlength{\ealglossraise}
\setlength{\ealglossraise}{1.15em}

% a tier 3 word, in English and in translation alike
\newcommand{\ealkey}[1]{\textcolor{ealkeycolour}{#1}}
\newcommand{\ealkeytr}[1]{\textcolor{ealkeycolour}{\ealtext{#1}}}

% One translated word above one English word. Built from kernel
% primitives, never a tabular, which would break wherever the sheet
% loads array, booktabs or tabularx. The box is as wide as the wider of
% the two, so a long translation pushes its neighbours aside instead of
% overlapping them, and it claims its full height so a table row or a
% TikZ node makes room for it.
\newcommand{\ealgloss}[2]{%
  \sbox{\ealboxen}{\ealkey{#1}}%
  \sbox{\ealboxtr}{\small\textcolor{ealtwocolour}{\ealtext{#2}}}%
  \setlength{\ealglosswd}{\wd\ealboxen}%
  \ifdim\wd\ealboxtr>\ealglosswd\setlength{\ealglosswd}{\wd\ealboxtr}\fi
  \leavevmode
  \makebox[\ealglosswd][c]{%
    \makebox[0pt][c]{\usebox{\ealboxen}}%
    \makebox[0pt][c]{\raisebox{\ealglossraise}{\usebox{\ealboxtr}}}%
  }%
}

% opens up line spacing around a block of glosses, so the raised
% translations sit clear of the line above
\newenvironment{ealglossed}{\par\linespread{2.1}\selectfont\ignorespaces}{\par}

% a whole translated sentence above its English counterpart. block level,
% so it does not belong inside a tabular cell
\newcommand{\ealpara}[2]{%
  \par\addvspace{0.5\baselineskip}%
  {\color{ealtwocolour}\ealtextblock{#1}}%
  \penalty200
  {\color{ealenglish}#2\par}%
  \addvspace{0.5\baselineskip}%
}

\begin{document}

% ---- Title ----
\begin{frame}[plain]
  \centering
  {\LARGE\bfseries \ealpara{\ealkeytr{تسلسل} کے اسٹارٹرز}{\ealkey{Sequences} Starters}}\\[0.5em]
\end{frame}

% ================================================================
% STARTER 1
% ================================================================
\begin{frame}[allowframebreaks]{Starter 1}

  \ealpara{Q1۔ ماچس کی تیلیوں کا نمونہ}{Q1.\enspace Matchstick Pattern}

  \matchsquares[red]{1}\quad
  \matchsquares[blue]{2}\quad
  \matchsquares[green!70!black]{3}\quad
  \ashowq{\matchsquares[orange]{4}}\\[0.4em]

  \ealpara{(a) چوتھی شکل بنائیں۔}{(a)\;Draw the 4th shape.}

  \ealpara{(b) یہ کتنی ماچس کی تیلیاں استعمال کرتی ہے؟}{(b)\;How many matchsticks does it use?}

  \ealpara{(b) 13 ماچس کی تیلیاں}{(b)\;13 matchsticks}

  \ealpara{Q2۔ \ealkeytr{تسلسل} مکمل کریں۔}{Q2.\enspace Complete the \ealkey{Sequence}}

  5,\;10,\;15,\;20,\;
  \ablank{25},\;\ablank{30},\;\ablank{35}

  \ealpara{(ہر بار 5 جمع کریں)}{(add 5 each time)}

  \ealpara{Q3۔ \ealkeytr{تسلسل} کے پہلے 5 \ealkeytr{ارکان} لکھیں جس کا \ealkeytr{پوزیشن ٹو ٹرم رول} 2n ہے۔}{Q3.\enspace Write the first 5 \ealkey{terms} of the \ealkey{sequence} with \ealkey{position-to-term rule}: $\boldsymbol{2n}$}

  \ablank{2},\;\ablank{4},\;\ablank{6},\;\ablank{8},\;\ablank{10}

  \ealpara{Q4۔}{Q4.\enspace}

  \ealpara{سیم 3 کا \ealkeytr{پہاڑا} لکھتا ہے: 3، 6، 9، 12، 15...}{Sam writes the 3 \ealkey{times table}: $3, 6, 9, 12, 15\ldots$}

  \ealpara{پھر وہ ہر \ealkeytr{رکن} میں 2 جمع کرتا ہے۔}{He then \textbf{adds 2} to every \ealkey{term}.}

  \ealpara{سیم کے \ealkeytr{تسلسل} کے پہلے 5 \ealkeytr{ارکان} لکھیں۔}{Write the first 5 \ealkey{terms} of Sam's \ealkey{sequence}.}

  \ealpara{\ealkeytr{عمومی رکن} کا اصول کیا ہے؟}{What is the \ealkey{nth term} rule?}

  \ealpara{\ealkeytr{ارکان}: 5، 8، 11، 14، 17۔ \ealkeytr{عمومی رکن}: 3n+2}{Terms:\;$5,\ 8,\ 11,\ 14,\ 17$.\quad \ealkey{nth term}:\;$\boldsymbol{3n+2}$}

  \ealpara{Q5۔ Q1 کو بڑھانا}{Q5.\enspace Extending Q1}

  \ealpara{Q1 کے ماچس کے نمونے کا استعمال کرتے ہوئے:}{Using the matchstick pattern from Q1:}

  \ealpara{(a) 10ویں شکل میں کتنی ماچس کی تیلیاں ہیں؟}{(a)\;How many matchsticks in the \textbf{10th} shape?}

  \ealpara{(b) nویں شکل کے لیے ایک اصول لکھیں۔}{(b)\;Write a rule for the nth shape.}

  \ealpara{(a) 31 ماچس کی تیلیاں۔}{(a)\;31 matchsticks.}

  \ealpara{(b) \ealkeytr{پوزیشن ٹو ٹرم رول}: 3n+1}{(b)\;\ealkey{position-to-term rule}:\;$3n+1$}

  \ealpara{Q6۔ چیلنج}{Q6.\;Challenge}

  \ealpara{کیا 100 اس \ealkeytr{تسلسل} کا ایک \ealkeytr{رکن} ہے جس کا \ealkeytr{عمومی رکن} 3n+1 ہے؟}{Is $\boldsymbol{100}$ a \ealkey{term} in the \ealkey{sequence} with \ealkey{nth term} $3n+1$\,?}

  \ealpara{دکھائیں کہ آپ کیسے جانتے ہیں۔}{Show how you know.}

  \ealpara{جی ہاں، جب n=33، 3n+1=3×33+1=99+1=100}{Yes, when $n=33$, $3n+1=3\times33+1=99+1=100$}

\end{frame}

% ================================================================
% STARTER 2
% ================================================================
\begin{frame}[allowframebreaks]{Starter 2}

  \ealpara{Q1۔ نقطوں کا نمونہ}{Q1.\enspace Dot Pattern}

  \fourtabledots[red]{1}\quad
  \fourtabledots[blue]{2}\quad
  \fourtabledots[green!70!black]{3}\quad
  \ashowq{\fourtabledots[orange]{4}}\\[0.4em]

  \ealpara{(a) چوتھی تصویر بنائیں۔}{(a)\;Draw the 4th picture.}

  \ealpara{(b) اس میں کتنے نقطے ہیں؟}{(b)\;How many dots does it contain?}

  \ealpara{(b) 16 نقطے}{(b)\;16 dots}

  \ealpara{Q2۔ \ealkeytr{تسلسل} مکمل کریں۔}{Q2.\enspace Complete the \ealkey{Sequence}}

  4,\;8,\;12,\;16,\;
  \ablank{20},\;\ablank{24},\;\ablank{28}

  \ealpara{(ہر بار 4 جمع کریں)}{(add 4 each time)}

  \ealpara{Q3۔ 3n+1 کے پہلے 5 \ealkeytr{ارکان} لکھیں۔}{Q3.\enspace Write the first 5 \ealkey{terms} of $\boldsymbol{3n+1}$}

  \ablank{4},\;\ablank{7},\;\ablank{10},\;\ablank{13},\;\ablank{16}

  \ealpara{Q4۔}{Q4.\enspace}

  \ealpara{ایلکس 5 کا \ealkeytr{پہاڑا} لکھتا ہے: 5، 10، 15، 20، 25...}{Alex writes the 5 \ealkey{times table}: $5, 10, 15, 20, 25\ldots$}

  \ealpara{پھر وہ ہر \ealkeytr{رکن} میں سے 1 تفریق کرتا ہے۔}{He then \textbf{subtracts 1} from every \ealkey{term}.}

  \ealpara{ایلکس کے \ealkeytr{تسلسل} کے پہلے 5 \ealkeytr{ارکان} لکھیں۔}{Write the first 5 \ealkey{terms} of Alex's \ealkey{sequence}.}

  \ealpara{\ealkeytr{عمومی رکن} کا اصول کیا ہے؟}{What is the \ealkey{nth term} rule?}

  \ealpara{\ealkeytr{ارکان}: 4، 9، 14، 19، 24۔ \ealkeytr{عمومی رکن}: 5n-1}{Terms:\;$4,\ 9,\ 14,\ 19,\ 24$.\quad \ealkey{nth term}:\;$\boldsymbol{5n-1}$}

  \ealpara{Q5۔ Q1 کو بڑھائیں}{Q5.\enspace Extend Q1}

  \ealpara{Q1 کے اپنے نقطوں کے نمونے کا استعمال کرتے ہوئے:}{Using your dot pattern from Q1:}

  \ealpara{(a) 10ویں تصویر میں کتنے نقطے ہیں؟}{(a)\;How many dots in the \textbf{10th} picture?}

  \ealpara{(b) nویں تصویر کے لیے ایک اصول لکھیں۔}{(b)\;Write a rule for the nth picture.}

  \ealpara{(a) 40 نقطے۔ (b) اصول: 4n}{(a)\;40 dots.\quad (b)\;Rule:\;$4n$}

  \ealpara{Q6۔ چیلنج}{Q6.\;Challenge}

  \ealpara{\ealkeytr{تسلسل} A کا \ealkeytr{عمومی رکن} 4n-2 ہے۔ \ealkeytr{تسلسل} B کا \ealkeytr{عمومی رکن} 2n+6 ہے۔ کیا وہ ایک ہی پوزیشن پر کوئی \ealkeytr{رکن} مشترک رکھتے ہیں؟}{\ealkey{Sequence} A has \ealkey{nth term} $4n-2$.\quad \ealkey{Sequence} B has \ealkey{nth term} $2n+6$.\quad Do they share a \ealkey{term} \emph{at the same position}?}

  \ealpara{4n-2=2n+6 ⇒ n=4۔ دونوں پوزیشن n=4 پر 14 کے برابر ہیں۔}{$4n-2=2n+6\;\Rightarrow\;n=4.$\; Both equal\;\textbf{14} at position $n=4$.}

\end{frame}

% ================================================================
% STARTER 3
% ================================================================
\begin{frame}[allowframebreaks]{Starter 3}

  \ealpara{Q1۔ ماچس کی تیلیوں سے مثلث}{Q1.\enspace Matchstick Triangles}

  \matchtriangles[red]{1}\quad
  \matchtriangles[blue]{2}\quad
  \matchtriangles[green!70!black]{3}\quad
  \ashowq{\matchtriangles[orange]{4}}\\[0.4em]

  \ealpara{(a) چوتھی شکل بنائیں۔}{(a)\;Draw the 4th shape.}

  \ealpara{(b) یہ کتنی ماچس کی تیلیاں استعمال کرتی ہے؟}{(b)\;How many matchsticks does it use?}

  \ealpara{(b) 12 ماچس کی تیلیاں}{(b)\;12 matchsticks}

  \ealpara{Q2۔ \ealkeytr{تسلسل} مکمل کریں۔}{Q2.\enspace Complete the \ealkey{Sequence}}

  7,\;14,\;21,\;28,\;
  \ablank{35},\;\ablank{42},\;\ablank{49}

  \ealpara{(ہر بار 7 جمع کریں --- 7 کا \ealkeytr{پہاڑا})}{(add 7 each time --- the 7$\times$ table)}

  \ealpara{Q3۔ 5n-2 کے پہلے 5 \ealkeytr{ارکان} لکھیں۔}{Q3.\enspace Write the first 5 \ealkey{terms} of $\boldsymbol{5n-2}$}

  \ablank{3},\;\ablank{8},\;\ablank{13},\;\ablank{18},\;\ablank{23}

  \ealpara{Q4۔}{Q4.\enspace}

  \ealpara{ایک \ealkeytr{تسلسل} 5 سے شروع ہوتا ہے اور ہر بار 3 بڑھتا ہے۔ پہلے 5 \ealkeytr{ارکان} لکھیں۔ \ealkeytr{عمومی رکن} کا اصول کیا ہے؟}{A \ealkey{sequence} \textbf{starts at 5} and \textbf{increases by 3} each time.\quad Write the first 5 \ealkey{terms}.\quad What is the \ealkey{nth term} rule?}

  \ealpara{\ealkeytr{ارکان}: 5، 8، 11، 14، 17۔ \ealkeytr{عمومی رکن}: 3n+2}{Terms:\;$5,\ 8,\ 11,\ 14,\ 17$.\quad \ealkey{nth term}:\;$\boldsymbol{3n+2}$}

  \ealpara{Q5۔ Q1 کو بڑھانا}{Q5.\enspace Extending Q1}

  \ealpara{Q1 کے اپنے مثلث کے نمونے کا استعمال کرتے ہوئے:}{Using your triangle pattern from Q1:}

  \ealpara{(a) nویں شکل کے لیے ایک اصول لکھیں۔}{(a)\;Write a rule for the nth shape.}

  \ealpara{(b) 20ویں شکل میں کتنی ماچس کی تیلیاں ہیں؟}{(b)\;How many matchsticks in the \textbf{20th} shape?}

  \ealpara{(a) اصول: 3n۔ (b) 60 ماچس کی تیلیاں۔}{(a)\;Rule:\;$3n$.\quad (b)\;60 matchsticks.}

  \ealpara{Q6۔ چیلنج}{Q6.\;Challenge}

  \ealpara{ایک \ealkeytr{تسلسل} کا \ealkeytr{عمومی رکن} 4n+3 ہے۔ کون سا \ealkeytr{رکن} پہلا ہے جو 75 سے تجاوز کرتا ہے؟ اپنا کام صاف صاف دکھائیں۔}{The \ealkey{nth term} of a \ealkey{sequence} is $4n+3$.\quad Which \ealkey{term} is the \textbf{first to exceed 75}?\quad Show your working clearly.}

  \ealpara{4n+3>75 ⇒ n>18۔ 19واں \ealkeytr{رکن}: 4(19)+3=79۔}{$4n+3>75\;\Rightarrow\;n>18.$\; 19th \ealkey{term}:\;$4(19)+3=\boldsymbol{79}.$}

\end{frame}

% ================================================================
% STARTER 4
% ================================================================
\begin{frame}[allowframebreaks]{Starter 4}

  \ealpara{Q1۔ نقطے}{Q1.\enspace Dots}

  \rectdots[red]{1}\quad
  \rectdots[blue]{2}\quad
  \rectdots[green!70!black]{3}\quad
  \ashowq{\rectdots[orange]{4}}\\[0.4em]

  \ealpara{(a) چوتھا نمونہ بنائیں۔}{(a)\;Draw the 4th pattern.}

  \ealpara{(b) اس میں کتنے نقطے ہیں؟}{(b)\;How many dots does it contain?}

  \ealpara{(b) 20 نقطے}{(b)\;20 dots}

  \ealpara{Q2۔ \ealkeytr{تسلسل} مکمل کریں۔}{Q2.\enspace Complete the \ealkey{Sequence}}

  1,\;5,\;9,\;13,\;
  \ablank{17},\;\ablank{21},\;\ablank{25}

  \ealpara{(ہر بار 4 جمع کریں)}{(add 4 each time)}

  \ealpara{Q3۔ 10-2n کے پہلے 5 \ealkeytr{ارکان} لکھیں۔}{Q3.\enspace Write the first 5 \ealkey{terms} of $\boldsymbol{10-2n}$}

  \ablank{8},\;\ablank{6},\;\ablank{4},\;\ablank{2},\;\ablank{0}

  \ealpara{Q4۔}{Q4.\enspace}

  \ealpara{میا 6 کا \ealkeytr{پہاڑا} لکھتی ہے: 6، 12، 18، 24، 30... پھر وہ ہر \ealkeytr{رکن} میں سے 4 تفریق کرتی ہے۔ پہلے 5 \ealkeytr{ارکان} لکھیں، \ealkeytr{عمومی رکن} معلوم کریں، اور 50واں \ealkeytr{رکن} معلوم کریں۔}{Mia writes the 6 \ealkey{times table}: $6, 12, 18, 24, 30\ldots$\quad She then \textbf{subtracts 4} from every \ealkey{term}.\quad Write the first 5 \ealkey{terms}, find the \ealkey{nth term}, and find the 50th \ealkey{term}.}

  \ealpara{\ealkeytr{ارکان}: 2، 8، 14، 20، 26۔ \ealkeytr{عمومی رکن}: 6n-4۔ 50واں \ealkeytr{رکن}: 296۔}{Terms:\;$2,\ 8,\ 14,\ 20,\ 26$.\quad \ealkey{nth term}:\;$6n-4$.\quad 50th \ealkey{term}:\;$\boldsymbol{296}$.}

  \ealpara{Q5۔ \ealkeytr{عمومی رکن} معلوم کریں}{Q5.\enspace Find the \ealkey{nth Term}}

  \ealpara{1، 5، 9، 13، 17، ... \ealkeytr{عمومی رکن} کا اصول کیا ہے؟}{$1,\;5,\;9,\;13,\;17,\;\ldots$\quad What is the \ealkey{nth term} rule?}

  \ealpara{\ealkeytr{فرق}: +4۔ پہلا \ealkeytr{رکن}: 1۔ اصول: 4n-3}{\ealkey{differences}: $+4$. \;First \ealkey{term}: 1.\quad Rule:\;$\boldsymbol{4n-3}$}

  \ealpara{Q6۔ چیلنج}{Q6.\;Challenge}

  \ealpara{ایک \ealkeytr{تسلسل} کا \ealkeytr{عمومی رکن} 10-2n ہے۔ آخری \ealkeytr{مثبت} \ealkeytr{رکن} کون سا ہے؟ پہلا \ealkeytr{منفی} \ealkeytr{رکن} کیا ہے؟}{The \ealkey{nth term} of a \ealkey{sequence} is $10-2n$.\quad Which is the \textbf{last \ealkey{positive} \ealkey{term}}?\quad What is the \textbf{first \ealkey{negative} \ealkey{term}}?}

  \ealpara{n=4: +2 (آخری \ealkeytr{مثبت})۔ n=5: 0۔ n=6: -2 (پہلا \ealkeytr{منفی})۔}{$n=4$:\;$+2$ (last \ealkey{positive}).\quad $n=5$:\;$0$.\quad $n=6$:\;$-2$ (first \ealkey{negative}).}

\end{frame}

% ================================================================
% STARTER 5
% ================================================================
\begin{frame}[allowframebreaks]{Starter 5}

  \ealpara{Q1۔ نقطے}{Q1.\enspace Dots}

  \crossdots[red]{1}\quad
  \crossdots[blue]{2}\quad
  \crossdots[green!70!black]{3}\quad
  \ashowq{\crossdots[orange]{4}}\\[0.4em]

  \ealpara{(a) چوتھی شکل بنائیں۔}{(a)\;Draw the 4th shape.}

  \ealpara{(b) شکل 5 میں کتنے نقطے ہیں؟}{(b)\;How many dots in shape 5?}

  \ealpara{(b) شکل 5 میں 17 نقطے ہیں}{(b)\;Shape 5 has 17 dots}

  \ealpara{Q5۔}{Q5.\enspace}

  \ealpara{جیک کہتا ہے: "میں 3 سے شروع کرتا ہوں اور ہر بار 4 جمع کرتا ہوں۔" (a) اس کے پہلے 5 \ealkeytr{ارکان} لکھیں، (b) \ealkeytr{عمومی رکن} کا اصول معلوم کریں، اور (c) فیصلہ کریں کہ کیا 59 اس کے \ealkeytr{تسلسل} میں ہے۔ اپنی استدلال دکھائیں۔}{Jake says: \\ \textit{``I start at 3 and add 4 each time.''} \\ (a) Write his first 5 \ealkey{terms}, (b) find the \ealkey{nth term} rule, and (c) decide whether $\boldsymbol{59}$ is in his \ealkey{sequence}. Show your reasoning.}

  \ealpara{(a) 3،7،11،15،19۔ (b) \ealkeytr{عمومی رکن}: 4n-1۔ (c) 4n-1=59 ⇒ n=15۔ جی ہاں! (15واں \ealkeytr{رکن})}{(a) $3,7,11,15,19$.\quad (b) \ealkey{nth term}:\;$4n-1$.\quad (c) $4n-1=59\;\Rightarrow\;n=15.$\;\textbf{Yes!} (15th \ealkey{term})}

  \ealpara{Q2۔ \ealkeytr{تسلسل} مکمل کریں۔}{Q2.\enspace Complete the \ealkey{Sequence}}

  1,\;4,\;9,\;16,\;
  \ablank{25},\;\ablank{36},\;\ablank{49}

  \ealpara{(\ealkeytr{مربع عدد}، n^2)}{(\ealkey{square numbers},\;$n^2$)}

  \ealpara{Q3۔ 2n+5 کے پہلے 5 \ealkeytr{ارکان} لکھیں۔}{Q3.\enspace Write the first 5 \ealkey{terms} of $\boldsymbol{2n+5}$}

  \ablank{7},\;\ablank{9},\;\ablank{11},\;\ablank{13},\;\ablank{15}

  \ealpara{Q4۔ \ealkeytr{عمومی رکن} معلوم کریں}{Q4.\enspace Find the \ealkey{nth Term}}

  \ealpara{(a) 3، 5، 7، 9، 11، ...  2n+1 (b) 6، 10، 14، 18، 22، ...  4n+2}{(a)\;$3,\;5,\;7,\;9,\;11,\;\ldots$\quad $2n+1$\quad (b)\;$6,\;10,\;14,\;18,\;22,\;\ldots$\quad $4n+2$}

  \ealpara{Q6۔ چیلنج}{Q6.\;Challenge}

  \ealpara{میا کے \ealkeytr{تسلسل} کا \ealkeytr{عمومی رکن} 3n+1 ہے۔ جیک کے \ealkeytr{تسلسل} کا \ealkeytr{عمومی رکن} 2n+6 ہے۔ پہلی تین قدریں معلوم کریں جو دونوں \ealkeytr{تسلسل} میں آتی ہیں۔ ان کے بارے میں آپ کیا نوٹ کرتے ہیں؟}{Mia's \ealkey{sequence} has \ealkey{nth term} $3n+1$.\quad Jake's \ealkey{sequence} has \ealkey{nth term} $2n+6$.\quad Find the \textbf{first three values} that appear in \emph{both} \ealkey{sequence}s. What do you notice about them?}

  \ealpara{مشترک قدریں: 10، 16، 22، ... (وہ ہر بار 6 بڑھتی ہیں)۔}{Shared values:\;$10,\;16,\;22,\;\ldots$\; (they increase by 6 each time).}

\end{frame}

% ================================================================
% STARTER 6
% ================================================================
\begin{frame}[allowframebreaks]{Starter 6}

  \ealpara{Q1۔ ماچس کی تیلیوں سے مثلث}{Q1.\enspace Matchstick Triangles}

  \matchtriangles[red]{1}\quad
  \matchtriangles[blue]{2}\quad
  \matchtriangles[green!70!black]{3}\quad
  \ashowq{\matchtriangles[orange]{4}}\\[0.4em]

  \ealpara{(a) چوتھی شکل بنائیں۔}{(a)\;Draw the 4th shape.}

  \ealpara{(b) یہ کتنی ماچس کی تیلیاں استعمال کرتی ہے؟}{(b)\;How many matchsticks does it use?}

  \ealpara{(c) یہ \ealkeytr{تسلسل} کون سا \ealkeytr{پہاڑا} ظاہر کرتا ہے؟}{(c)\;Which \ealkey{times tables} does the \ealkey{sequence} represent?}

  \ealpara{(b) 12 ماچس کی تیلیاں (c) 3 کا \ealkeytr{پہاڑا}}{(b)\;12 matchsticks\quad (c)\;The 3 \ealkey{times table}}

  \ealpara{Q2۔ \ealkeytr{تسلسل} مکمل کریں۔}{Q2.\enspace Complete the \ealkey{Sequence}}

  4,\;8,\;12,\;16,\;
  \ablank{20},\;\ablank{24},\;\ablank{28}

  \ealpara{Q3۔ 3 کا \ealkeytr{پہاڑا} کا \ealkeytr{عمومی رکن} کا اصول 3n ہے، 5 کا \ealkeytr{پہاڑا} کا \ealkeytr{عمومی رکن} کا اصول کیا ہے؟}{Q3.\enspace The 3 \ealkey{times table} has \ealkey{nth term} rule $3n$, what is the \ealkey{nth term} rule for the 5 \ealkey{times table}?}

  \ablank{$5n$}

  \ealpara{Q4۔ 3n+4 کے پہلے 5 \ealkeytr{ارکان} لکھیں۔}{Q4.\enspace Write the first 5 \ealkey{terms} of $\boldsymbol{3n+4}$}

  \ablank{7},\;\ablank{10},\;\ablank{13},\;\ablank{16},\;\ablank{19}

  \ealpara{Q5۔}{Q5.\enspace}

  \ealpara{ایک \ealkeytr{تسلسل} 3 کا \ealkeytr{پہاڑا} لے کر اور ہر \ealkeytr{رکن} میں سے 2 تفریق کر کے بنایا گیا ہے۔ (a) پہلے 5 \ealkeytr{ارکان} لکھیں۔ (b) \ealkeytr{عمومی رکن} کا اصول کیا ہے؟ (c) اپنے اصول کا استعمال کرتے ہوئے 20واں \ealkeytr{رکن} معلوم کریں۔ (d) کیا 100 اس \ealkeytr{تسلسل} میں ہے؟ اپنا کام صاف صاف دکھائیں۔}{A \ealkey{sequence} is made by taking the 3 \ealkey{times table} and \textbf{subtracting 2 from every \ealkey{term}.}\quad (a)\;Write the first 5 \ealkey{terms}.\quad (b)\;What is the \ealkey{nth term} rule?\quad (c)\;Use your rule to find the \textbf{20th} \ealkey{term}.\quad (d)\;Is $\boldsymbol{100}$ in the \ealkey{sequence}? Show your working clearly.}

  \ealpara{(a) 1،4،7،10،13۔ (b) 3n-2۔ (c) 3×20-2=58۔ (d) 3n-2=100 ⇒ 3n=102 ⇒ n=34۔ جی ہاں، 100 34واں \ealkeytr{رکن} ہے۔}{(a)\;1,\,4,\,7,\,10,\,13. \quad (b)\;$\boldsymbol{3n-2}$. \quad (c)\;$3\times20-2 = 58$.\quad (d)\;$3n-2=100 \Rightarrow 3n=102 \Rightarrow n=34$. \textbf{Yes}, 100 is the 34th \ealkey{term}.}

\end{frame}

% ================================================================
% STARTER 7
% ================================================================
\begin{frame}[allowframebreaks]{Starter 7}

  \ealpara{Q1۔ نقطے}{Q1.\enspace Dots}

  \fourtabledots[red]{1}\quad
  \fourtabledots[blue]{2}\quad
  \fourtabledots[green!70!black]{3}\quad
  \ashowq{\fourtabledots[orange]{4}}\\[0.4em]

  \ealpara{(a) چوتھی تصویر بنائیں۔}{(a)\;Draw the 4th picture.}

  \ealpara{(b) اس میں کتنے نقطے ہیں؟}{(b)\;How many dots does it contain?}

  \ealpara{(c) یہ \ealkeytr{تسلسل} کون سا \ealkeytr{پہاڑا} ظاہر کرتا ہے؟}{(c)\;Which \ealkey{times table} does the \ealkey{sequence} represent?}

  \ealpara{(b) 16 نقطے (c) 4 کا \ealkeytr{پہاڑا}}{(b)\;16 dots\quad (c)\;The 4 \ealkey{times table}}

  \ealpara{Q2۔ \ealkeytr{تسلسل} مکمل کریں۔}{Q2.\enspace Complete the \ealkey{Sequence}}

  5,\;10,\;15,\;20,\;
  \ablank{25},\;\ablank{30},\;\ablank{35}

  \ealpara{(ہر بار 5 جمع کریں --- 5 کا \ealkeytr{پہاڑا})}{(add 5 each time --- the 5 \ealkey{times table})}

  \ealpara{Q3۔ 4n-3 کے پہلے 5 \ealkeytr{ارکان} لکھیں۔}{Q3.\enspace Write the first 5 \ealkey{terms} of $\boldsymbol{4n-3}$}

  \ablank{1},\;\ablank{5},\;\ablank{9},\;\ablank{13},\;\ablank{17}

  \ealpara{Q4۔}{Q4.\enspace}

  \ealpara{5 کا \ealkeytr{پہاڑا} لے کر شروع کریں اور ہر \ealkeytr{رکن} میں 1 جمع کریں۔ (a) پہلے 5 \ealkeytr{ارکان} لکھیں۔ (b) \ealkeytr{عمومی رکن} کا اصول کیا ہے؟ (c) 20واں \ealkeytr{رکن} معلوم کریں۔ (d) کیا 104 اس \ealkeytr{تسلسل} میں ہے؟ وضاحت کریں کہ آپ کیسے جانتے ہیں۔}{Start with the \textbf{5 \ealkey{times table}} and \textbf{add 1} to every \ealkey{term}.\quad (a)\;Write the first 5 \ealkey{terms}.\quad (b)\;What is the \ealkey{nth term} rule?\quad (c)\;Find the 20th \ealkey{term}.\quad (d)\;Is $\boldsymbol{104}$ in the \ealkey{sequence}? Explain how you know.}

  \ealpara{(a) 6،11،16،21،26۔ (b) 5n+1۔ (c) 5×20+1=101۔ (d) 5n+1=104 ⇒ 5n=103 ⇒ n=20.6۔ نہیں، 104 \ealkeytr{رکن} نہیں ہے۔}{(a)\;6,\,11,\,16,\,21,\,26. \quad (b)\;$\boldsymbol{5n+1}$. \quad (c)\;$5\times20+1 = 101$.\quad (d)\;$5n+1=104 \Rightarrow 5n=103 \Rightarrow n=20.6$. \textbf{No}, 104 is not a \ealkey{term}.}

\end{frame}

\end{document}